\section{The finite rank-three host and exact automorphic endpoints}
\label{sec:sr-rank-three-host-endpoints}

\subsection{Constant characteristic polynomials and the Gaussian finite part}
\label{subsec:sr-constant-character-finite-part}
We make explicit the last algebraic step from
\(\mathrm{FinEnd}\) to a finite weight-zero algebra.
All groups below are the actual finitely ramified arithmetic
quotients, so Lemma~\ref{lem:sr-finite-cross-cocycles}
and Proposition~\ref{prop:sr-finite-CH} apply.
The latter proposition and the integral determinant identity
(H9) are valid in characteristic \(3\).

\begin{srlemma}\label{lem:sr-constant-character-triangular}
Let \(k\) be a finite field of odd characteristic,
\(A\) a local domain containing \(k\), and
\(r:\Gamma\to\mathrm{GL}_3(A)\rtimes\{\pm1\}\)
an actual normalized rank-three unitary point.
Write \(H=\ker\eta\).
Suppose its reduction is the torus tuple
\((\chi,1,\chi^{-1})\), with \(\chi(g_1)^2\ne1\).
Suppose on an open subgroup of \(H\) its matrix
characteristic polynomials equal those of that tuple.
Then the matrices on \(H\) are all upper triangular
or all lower triangular in the marked basis, and every
diagonal entry equals its prescribed constant character.
\end{srlemma}
\begin{proof}
For each \(g\in H\), some positive power \(g^s\)
belongs to the given open subgroup.
An eigenvalue \(\alpha\) of \(r(g)\) has \(\alpha^s\)
algebraic over \(k\), so \(\alpha\) is algebraic over \(k\).
Thus every characteristic coefficient of \(r(g)\)
is algebraic over \(k\).
An element of \(A\) algebraic over the finite field \(k\)
belongs to a finite subfield of \(A\).
Reduction is injective on that subfield.
Its prescribed residual value consequently determines it
exactly.
It follows that
\[
 \operatorname{char}_{r(g)}(T)
                   =(T-\chi(g))(T-1)(T-\chi(g)^{-1})
                                \quad(g\in H).
 \tag{CF1}
\]

Here is the exact triangularity argument in the marked basis.
Let \(c\) be the normalized conjugation, so its matrix
part is \(1\), and put \(g^*=cg^{-1}c\).
The unitary law gives
\[
 r(g^*)=J_3r(g)^{\mathsf t}J_3^{-1}\qquad(g\in H).
\]
The marked normalization kills the two adjacent entries
of \(r(g_1)-r(g_1^*)\).
The other two off-diagonal entries cancel automatically
under this transpose reflection.
Hence
\[
             r(g_1)-r(g_1^*)=\operatorname{diag}(\zeta,0,-\zeta),
 \tag{CF2}
\]
where \(\zeta\) reduces to
\(\chi(g_1)-\chi(g_1)^{-1}\) and is a unit.
The three diagonal values in (CF2) are distinct,
since \(2\ne0\).

We first exclude irreducibility over \(\operatorname{Frac}A\).
By (H9), the characteristic polynomials on group words
determine the full determinant on the group algebra
without division by \(3!\).
Equation (CF1) therefore factors this representation
through the constant-character Cayley--Hamilton algebra
of Proposition~\ref{prop:sr-finite-CH}, after scalar extension.
Indeed continuity modulo every power of the maximal
ideal gives a finite quotient representation at each
precision; its determinant kills the defining
Cayley--Hamilton relations, and separatedness then makes
the original representation kill their closed ideal.
Thus passage to the fraction field does not silently
discard the completed-algebra condition.
Its nilpotent radical has quotient \(k^3\), as in that
proposition's proof.
Every irreducible representation of this algebra is
one-dimensional: the radical annihilates a simple module,
and the three coordinate idempotents of \(k^3\) split it.
Thus the rank-three representation is reducible.

Any invariant subspace is invariant under (CF2).
Its spectral projectors, obtained by Lagrange interpolation
at \(\zeta,0,-\zeta\), show that it is spanned by
a nonempty proper subset of the three coordinate vectors.
If every matrix entry \(ij\) vanishes on \(H\), so does
entry \(2-j,2-i\), by the unitary reflection.
There are six coordinate subsets to check.
The subsets \(\{0\},\{0,1\}\) force upper triangularity;
\(\{2\},\{1,2\}\) force lower triangularity.
The subsets \(\{1\},\{0,2\}\) force a block decomposition
of the middle line and the outer two coordinates.

In this last case the middle entry is an eigenvalue,
hence algebraic over \(k\); its residual value is \(1\),
so it equals \(1\).
The outer block has characteristic polynomial
\((T-\chi(g))(T-\chi(g)^{-1})\) for every \(g\).
Apply the rank-two instance of the same
Cayley--Hamilton argument.
The outer representation is reducible.
Its invariant line is a coordinate line because
the restriction of (CF2) is
\(\operatorname{diag}(\zeta,-\zeta)\).
One of the two corner entries consequently vanishes
on every \(g\), and the whole rank-three representation
is upper or lower triangular.
Injection into \(\operatorname{Frac}A\) transports
these zero entries back to \(A\).

Each diagonal entry in a triangular matrix is an
eigenvalue, so the finite-field argument following
(CF1) fixes it to its prescribed residual character.
Finally the matrix part of \(r(gc)\) equals that of
\(r(g)\); thus the same diagonal assertion holds
outside \(H\) as well.
\end{proof}

\begin{srlemma}\label{lem:sr-graded-finite-part-criterion}
Let \(\mathcal O\) be a complete coefficient DVR with
finite residue field and uniformizer \(\varpi\).
Let \(R\) be the actual complete Noetherian local
homogeneously presented algebra, with finitely many
homogeneous generators \(z_i\in\mathfrak m_R\),
of integer weights.
Let \(B=R_0\) be its completed weight-zero algebra.
Suppose every local specialization \(R\to A\), with
\(A\) a complete characteristic-residue-prime DVR
with finite residue field, kills every zero-weight
generator and kills either all positive-weight or
all negative-weight generators.
Then \(B\) is finite over \(\mathcal O\).
\end{srlemma}
\begin{proof}
The zero-weight monoid of exponent vectors has finitely
many generators.
For example intersect the nonnegative rational cone
with the integral kernel of the weight homomorphism;
triangulating its rational polyhedral cones and taking
the finitely many lattice points in their bounded
fundamental parallelepipeds proves this assertion.
Consequently \(B\) is a complete Noetherian local
\(\mathcal O\)-algebra with maximal ideal generated by
\(\varpi\) and finitely many nonconstant balanced
monomials \(b_1,\ldots,b_t\) in the \(z_i\).
The continuous weight-zero projection is
\(\mathcal O\)-linear, and gives
\[
                       (\varpi R)\cap B=\varpi B.
 \tag{CF3}
\]

We claim \(b_j\in\sqrt{\varpi R}\).
Otherwise choose a prime \(P\supset(\varpi)\) avoiding
\(b_j\).
In the complete local domain \(R/P\), its image is
nonzero and belongs to the maximal ideal.
There is a one-dimensional domain quotient preserving
that nonzero image.
Here is a direct curve-selection argument.
If the dimension exceeds one, choose an element of
the maximal ideal outside the finitely many height-one
minimal primes above that image.
A minimal prime above the chosen element has height
one by the principal ideal theorem and avoids the
image; quotient by it.
Dimension decreases, and iteration ends in dimension one.
The normalization of this one-dimensional complete local
domain is finite, local and a complete DVR, with finite
residue field.
Thus there is a local map \(R\to A\) of the kind in
the hypothesis with \(b_j\) nonzero.

A nonconstant zero-weight monomial either contains a
zero-weight variable, or contains variables of both
signs.
Every such monomial maps to zero under the hypothesis,
a contradiction.
This proves the claim.
Use (CF3) on a power of \(b_j\) to get
\(b_j\in\sqrt{\varpi B}\).
Since these finitely many monomials and \(\varpi\)
generate the maximal ideal,
\(\mathfrak m_B^s\subset\varpi B\) for some \(s\).
Hence \(B/\varpi B\) is finite over the residue field.
The \(\varpi\)-adic and maximal-ideal topologies are
equivalent by this containment.
Completeness and the usual successive lifting of a
basis of \(B/\varpi B\) show that finitely many elements
generate \(B\) over \(\mathcal O\).
\end{proof}

\begin{srproposition}\label{prop:sr-host-finite-part-from-finend}
For the actual marked rank-three host algebra,
\(\mathrm{FinEnd}\) implies finiteness of its weight-zero
part over the coefficient DVR.
\end{srproposition}
\begin{proof}
Its adapted homogeneous generators are matrix entries
of \(r-r_{\mathrm{tor}}\); their weights are the
differences of \(1,0,-1\).
They all belong to the maximal ideal by residual
normalization.
A characteristic-\(\ell\) DVR specialization gives the
ordinary point to which \(\mathrm{FinEnd}\) applies.
That condition fixes its characteristic polynomials
on an open subgroup of \(\ker\eta\).
Lemma~\ref{lem:sr-constant-character-triangular} makes
it upper or lower triangular with the prescribed
constant diagonal.
The zero-weight generator images are therefore zero;
in one triangular case all positive-weight images
vanish, and in the other all negative-weight images
vanish.
Lemma~\ref{lem:sr-graded-finite-part-criterion}
gives the required finite weight-zero part.
\end{proof}

The exact frozen-source chain is
\srcid{CurveDVRSpecialization/ConstOfChar.lean},
\srcid{UnitaryThreeDeformations/UniversalDiagonalOfTriangularPoints.lean},
and
\srcid{CurveDVRSpecialization/FinCrit.lean}.
The preceding proofs give their needed arithmetic specialization.
The intermediate \srcid{DiagEq} theorem fixes diagonals for
general triangular ordinary points using the frame's class-one
condition; here characteristic-\(\ell\) finiteness needs only
the stronger constant-character argument (CF1).
No unspecified triangularity or diagonal-finiteness assertion
is retained in the implication from (G10) to the finite
Gaussian host algebra.


\subsection{The actual Gaussian seeds and their common support}
\label{subsec:sr-gaussian-producer}
We give the arithmetic specialization needed for the
large-conductor Gaussian branch, including coefficient
prime \(3\).
The shared sharp/flat propagation theorem has a precise
remaining role at the end of this subsection;
automorphy is not assumed as a Gaussian literature input.

Fix \(K=\mathbb Q(i)\), \(\eta=\eta_K\),
\(\ell\in\{3,7\}\), twist \(1\), and an actual principal
character \(\psi\) in the large-conductor branch.
Thus \(\ell\nmid N\), \(4\mid N\),
\(\psi(-1)=-1\), and its dyadic component \(\psi_2\)
is even and primitive of conductor \(2^c\).
Put \(\Psi=\epsilon_\ell\psi\).
The bounds in Lemma~\ref{lem:sr-dyadic-three-seven}
give \(c\ge14\) at \(3\), \(c\ge18\) at \(7\), and
\[
           \operatorname{ord}\psi(3)=2^{c-2},\qquad
           \operatorname{ord}\psi(7)=2^{c-3}.
 \tag{G1}
\]
In particular these orders exceed \(8\).
All finite character values have order prime to \(\ell\).

\subsubsection{A scalar prime and the real quadratic base}
Choose a dyadic unit \(b\equiv1\pmod4\) with
\(\psi_2(b)=-1\).
Such a unit exists because \(\psi_2\) is faithful on
the cyclic \(1+4\)-quotient.
Set the odd-prime residues equal to \(1\).
The Chinese remainder theorem and Dirichlet's theorem
give a fresh prime \(p_0\) with
\[
         p_0\equiv b\pmod{2^c},\quad
         p_0\equiv1\pmod{qd},\quad
         p_0\equiv-1\pmod\ell .
\]
The harmless larger dyadic level in \(N\) can be included
by lifting the chosen residue.
Hence
\[
       \eta(p_0)=1,\quad
       \bar\psi(p_0)=-1,\quad
       \bar\Psi(p_0)=1,\quad
       \bar\epsilon_\ell(p_0)=-1 .
 \tag{G2}
\]
Thus \(p_0\) splits in \(K\).
Both residual tuples
\((\bar\Psi,\bar\Psi^{-1})\) and
\((\bar\Psi,1,\bar\Psi^{-1})\) are scalar there.
We may avoid any given finite set.

Choose a fresh positive prime \(D\equiv7\pmod8\).
At each odd level prime \(p\), including \(\ell\),
require
\[
                   (D/p)=\eta(p),
\]
and also require \((D/p_0)=1\).
These conditions specify nonzero residues modulo the
distinct primes, so another reduced arithmetic progression
and Dirichlet's theorem produce such a \(D\).
Put
\[
                       F_0^+=\mathbb Q(\sqrt D).
\]
At \(2\), its completion is \(\mathbb Q_2(i)\);
at every odd level prime its quadratic local character
equals that of \(K/\mathbb Q\).
Thus every level prime splits in
\(F_0^+K/F_0^+\).
Moreover
\[
 \operatorname{disc}(F_0^+K)
        =(-4)(4D)(-D)=16D^2
                     =\operatorname{disc}(F_0^+)^2 .
\]
The relative discriminant formula shows that this CM
quadratic extension is unramified at every finite place.
This real quadratic field is the exact index-two
\(\mathrm{H}_0\) of the source's
\srcid{FrobeniusSplitPrimeDatum}.
The extra ramification at \(D\), outside the level,
supplies that datum's fresh ramified prime.
The prime \(p_0\) splits in \(F_0^+K\).
At \(\ell\), the completion of \(F_0^+\) is the
unramified quadratic extension, since \(\ell\) is
inert in \(K\).

Choose a fresh prime
\(w\equiv1\pmod{\ell N}\) and a character \(\vartheta\)
of exact order \(\ell\), conductor \(w\).
It is even and congruent to \(1\) modulo the coefficient
maximal ideal.
Let \(F_\vartheta/\mathbb Q\) be its real cyclic
degree-\(\ell\) field.
For the depth \(r\) required by the common frame put
\[
 F_r^+=F_0^+F_\vartheta
        \mathbb Q(\zeta_{\ell^{r+1}})^{\mu_{\ell-1}},
                  \qquad L_r=F_r^+K .
 \tag{G3}
\]
The CM extension remains finite-unramified and split
at all level primes.
Its real degree is even.
The new residue degrees at \(p_0\) are powers of the
odd prime \(\ell\), so (G2) survives on its decomposition
groups over \(F_r^+\).
This supplies the actual scalar-place condition for
both ranks and their scalar twists.

\subsubsection{The exact local conditions of the Gaussian frame}
There is no need to require
\(\bar\Psi^2\ne1\) on \(\ell\)-adic inertia.
That assertion would fail at \(\ell=3\).
The source's local separation condition concerns the
whole decomposition group, and it holds here.

Let \(f\) be a rational local uniformizer lift with
\(\epsilon_\ell(f)=1\).
Since \(K/\mathbb Q\) is inert at \(\ell\), \(f^2\)
fixes \(K\), and
\[
                 \bar\Psi(f^2)^2=\bar\psi(\ell)^4\ne1 .
 \tag{G4}
\]
Over a field in (G3) the unramified local degree is
\(2\ell^u\) for some \(u\ge0\);
the cyclotomic extra ramification has \(\ell\)-power degree.
Thus (G1) still gives (G4) on every actual local
decomposition group.
More explicitly, for \(j=\pm1,\pm2\), both
\(\bar\Psi^j\) and
\(\bar\Psi^j\bar\epsilon_\ell^{-1}\)
are nontrivial there.
Evaluate on \(f^{2\ell^u}\):
the cyclotomic value is \(1\) and the nontriviality
is detected by \(\bar\psi(\ell)^{2j\ell^u}\).
These are the full-local-group invariant vanishings
used in local Tate duality and the Euler formula.

The mod-\(\ell\) cyclotomic inertia image is still full:
the initial real base is unramified at \(\ell\) and
the later ramification has \(\ell\)-power index,
whereas its order is \(\ell-1\).
This is exactly the source's \(\mathrm{UFacts.mu}\).
The local fields contain no \(\mu_\ell\).
Their pro-\(\ell\) unit groups therefore have the
free ranks used in the ordinary module construction.
The character \(\bar\Psi\) is trivial on this pro-\(\ell\)
inertial part.

The other Gaussian frame separation has an explicit
global witness.
Take the dyadic residue \(-1\), odd-prime residues \(1\),
and residue \(1\pmod\ell\).
Evenness of \(\psi_2\) gives an element with
\[
              \bar\epsilon_\ell=1,\qquad
              \bar\Psi=1,\qquad \eta=-1 .
 \tag{G5}
\]
Thus
\(\ker\bar\epsilon_\ell\cap\ker\bar\Psi
 \not\subset\ker\eta\), the exact \(\mathrm{Gen.g4}\).
Equation (G4) supplies \(\mathrm{Gen.g5}\).
All finite inertial values away from \(\ell\) have
\(2\)-power order, so their Teichm\"uller lift is
unchanged and the source's character \(\mathrm{H}_b\)
may be the whole group.

Apply the shared arithmetic-frame construction with this
fixed real quadratic base and these witnesses.
Its Ferrero--Washington class-rank bound is chosen
before the cyclotomic depth; its local invariant
conditions are the full-decomposition-group assertions
just checked.
For the two strict classes
\(\bar\Psi^{-a'}\eta\), \(a'=1,2\), nontriviality over
\(F_r^+\) follows from (G4), including when inertia is
trivial at \(3\).
Consequently the restriction on \(H^1\) is injective
by inflation--restriction.
The local Euler/Poitou--Tate proof and the class-rank
bounds are unchanged.
Choose the depth only afterwards, as in that shared proof.
The marked element is chosen using (G4), in the local
decomposition group.
It is not required to belong to inertia.
These observations match the source's
\(\mathrm{Frame.Ok}\), \(\mathrm{Frame.Big}\) and
\(\mathrm{GalLink}\) conditions at \(3\) as well as \(7\).

\subsubsection{Two actual cuspidal seeds}
Use \(\chi'=\psi\), \(j=1\), and the same
\(\eta,\vartheta,w,L_r\) for both ranks.
These satisfy the source's
\srcid{FieldPairCharacterTripleConditions}:
\(\eta\) cuts out \(K\), its conductor divides \(N\),
\(\chi'(-1)=-1\), the prime \(w\) splits both finite
characters, and \(\vartheta|_{G_{L_r}}=1\).
The local witness (G4) is its nontrivial-square condition.
Also \(\chi'\ne\eta\) and its order is prime to \(\ell\),
by the large dyadic conductor.

Take \(m\equiv1\pmod{\ell-1}\) as large as needed.
For rank three, apply the constructed Eisenstein seed
to the primitive character underlying
\[
                         \varphi_3=\psi\eta
\]
and the auxiliary \(\vartheta\), at weight \(m+1\).
This even character is nontrivial: a conductor at least
\(2^{14}\) cannot be canceled by the quadratic
conductor-four factor.
Let \(\rho_3\) be the resulting two-dimensional
ordinary cuspidal Galois representation.
Set
\[
                            r_3=\operatorname{ad}^0\rho_3
                                                    |_{G_{L_r}} .
\]
Its residual tuple, ordinary inertial exponents and
finite inertia are exactly
\[
 \begin{array}{c|c|c}
 \hbox{residual tuple}&\hbox{cyclotomic exponents}&
                                     \hbox{finite inertia tuple}\\ \hline
 (\bar\Psi,1,\bar\Psi^{-1})&(m,0,-m)&(\psi,1,\psi^{-1}) .
 \end{array}
 \tag{G6}
\]
The scalar automorphic normalization is
\(k_3=\epsilon_\ell^{-1}\).

For rank two use instead
\(\varphi_2=\psi^2\eta\), an odd, hence nontrivial,
character, and classical weight \(2m+1\).
Let \(\rho_2\) be its ordinary cuspidal representation
and put
\[
                   \sigma_2=\rho_2\epsilon_\ell^{-m}\psi^{-1}.
\]
The determinant and ordered residual characters are
\[
 \det\sigma_2=\eta\vartheta,\qquad
 \bar\sigma_2^{\,\mathrm{ss}}
       =(\bar\epsilon_\ell^m\bar\psi\eta)
                     \oplus(\bar\epsilon_\ell^{-m}\bar\psi^{-1}).
 \tag{G7}
\]
Its restriction \(r_2=\sigma_2|_{G_{L_r}}\) therefore
has determinant one, residual tuple
\((\bar\Psi,\bar\Psi^{-1})\), ordinary exponents
\((m,-m)\), and finite inertia
\((\psi,\psi^{-1})\), with the indicated ordering.
These are the exact rank-two seed data,
not just a matching adjoint trace.

Both classical levels are their primitive character
conductors times \(w\), and divide \(Nw\).
Their minimal-level local calculation gives \(N=0\)
and the specified rank-one inertial summands.
At \(w\), restriction to \(L_r\) kills
\(\vartheta\), so both seeds become unramified there.
Thus \(w\) does not enlarge their tame packet on \(L_r\).
At \(\ell\) they are crystalline ordinary, with the
unit-root congruence supplied by the Eisenstein seed.
The coefficient-\(3\) construction is
Section~\ref{subsec:sr-eisenstein-three}; no
\(\ell\ge5\) assumption is retained.

The strong-open irreducibility proof in the constructed
seed applies to both \(\rho_i\).
For \(\rho_3\), the ratio of the residual characters
on a local element above \(\ell\) has nontrivial square
by (G1)--(G4).
For \(\rho_2\), its square is detected by
\(\psi(\ell)^8\ne1\).
Thus neither residual ratio can be quadratic.
A quadratic self-twist is excluded, and the distinct
Hodge numbers exclude finite projective image.
Consequently both two-dimensional representations and
the rank-three adjoint remain irreducible on every
open subgroup.
This proof is valid in coefficient characteristic
\(3\), since only characteristic zero and preservation
of quadratic characters modulo an odd prime are used.

\subsubsection{An explicit Gaussian rank-two scalar character}
There is a particularly concrete \(k_2\) in this case.
Let \(E:y^2=x^3-x\).
Its CM endomorphism \(i\) is defined over \(K\);
after splitting coefficients its Tate representation on
\(G_K\) is a sum of characters \(\lambda,\lambda^c\).
The Weil pairing gives
\[
                              \lambda\lambda^c=\epsilon_\ell .
\]
The character \(\lambda\) is the usual Gaussian
algebraic Hecke character: a prime-to-\(2\) ideal has
a unique generator congruent to \(1\) modulo
\((1+i)^3\), and this generator gives its ideal
character, with the reciprocity orientation fixed above.
This explicit description is recorded, for example, in
\href{https://www.college-de-france.fr/media/professeurs-honoraires/UPL6386850193290496047_cours_Taiwan_avril__10juin_2011__5.pdf}
{Serre, \emph{Lectures on \(N_X(p)\)}, Section 2.2.2}.
It shows that the finite conductor is supported at \(2\)
and its finite inertial values lie in \(\mu_4\).
Choose the CM embedding for which \(\lambda\) has
cyclotomic exponents \((1,0)\), and set
\[
                                  k_2=\lambda^{-1}|_{G_{L_r}} .
 \tag{G8}
\]
It is algebraic, has exponents \((-1,0)\), and
\(k_2k_2^c=\epsilon_\ell^{-1}\).

This character also satisfies the source's transfer
condition, which cannot be replaced just by that norm
identity.
For \(t\in G_{\mathbb Q}\setminus G_K\), the matrix of
\(t\) on the two CM lines is off-diagonal.
Its determinant is \(\epsilon_\ell(t)\).
Its square on either line is therefore
\(-\epsilon_\ell(t)\); hence
\[
              \lambda(t^2)=-\epsilon_\ell(t),\qquad
              \lambda^{-1}(t^2)=-\epsilon_\ell(t)^{-1}.
 \tag{G9}
\]
This is exactly \(\mathrm{IsTwist.ver}\).
Finite inertia away from \(\ell\) is killed by the
prime-to-\(\ell\) integer \(4\).
Good reduction of \(E\) at \(3,7\) makes its rank-one
characters crystalline there.
All fields of coefficients are finite.
Thus (G8) supplies every character requirement in
the actual rank-two seed normalization.

The trace-zero trace form gives the symmetric
conjugate pairing of \(r_3\).
For rank two, write \(J\) for the alternating
two-by-two form and use \(J\sigma_2(c)\).
Since \(\det\sigma_2(c)=-1\) and \(c^2=1\),
this matrix is symmetric and gives a conjugate pairing
of \(r_2\) with multiplier one.
Twisting by (G8) changes that multiplier to
\(\epsilon_\ell^{-1}\), the rank-two conjugate-self-dual
Galois normalization.
Similarly \(r_3\epsilon_\ell^{-1}\) has multiplier
\(\epsilon_\ell^{-2}\), the rank-three normalization.

Cyclic base change of the two classical cusp forms
to \(L_r\) is cuspidal by strong-open irreducibility.
Their algebraic twists remain cuspidal and regular.
The Gelbart--Jacquet adjoint lift gives the cuspidal
rank-three realization of (G6), with normalization
\(r_3\epsilon_\ell^{-1}\).
The actual local principal-series parameters have
\(N=0\).
The ordinary eigenvalues are those determined by
the unit and uniformizer equations (H5)--(H6), in
ranks two and three respectively.
The generic unitary window, or the adjoint root bound
proved in the constructed seed, supplies the
unlinked local condition.
This verifies the full automorphic-and-local seed
conditions \(\mathrm{IsSeedAlt}\).

\subsubsection{Insertion in the same two function modules}
The definite groups \(U(2)\) and \(U(3)\) may both be
taken from the positive definite identity hermitian forms.
At every finite place the CM extension is unramified,
and units are local norms; these forms are quasi-split
there.
Their real signatures are definite, and the even
real degree in (G3) gives the required existence parity.
They are actual pure inner forms.

Apply the stable cuspidal seed insertion argument of
Lemma~\ref{lem:sr-stable-three-seed-insertion} in ranks
three and two.
Only the rank of the local type and Jacquet equations
changes.
KMSW is used only for its proved generic-parameter
pure-inner case.
At tame places the coefficient lattice is
\(\bigotimes\sigma^{\mathrm{crys}}(\tau)^\vee\).
At \(\ell\) the ordered ordinary Jacquet vector has
exactly the \(\theta^{-1}\) coefficient character and
the normalized unit \(U_j\)-values in (H5)--(H6).
The finite integral vector, its residual localization
and classical control therefore occur in the same
ordinary modules and universal determinant quotients
used for the two runs.
The actual Hecke kernels vanish at these actual seeds.
This supplies the Gaussian \(\mathrm{seedPtOf}\)
hypotheses in both ranks.

The coefficient interpolation, ordinary projector,
projectivity, coefficient change, old/new maps and
nilpotence calculation of the common construction
are rank independent.
They divide neither by \(3\) nor by \(3!\);
the rank-three nilpotence exponent is \(3!=6\).
The normalized inertia types and ordinary characters
here have just been identified.
Accordingly the only shared support input still needed
at this interface is the completed sharp/flat
rank-two/rank-three propagation theorem, with its
actual arithmetic frame bound and its two strict
classes.
It is the theorem implemented by
\(\mathrm{PrLinkedEnd}\to\mathrm{AbsMain}
 \to\mathrm{PropagationTheoremGL3}\), using the
rank-two sharp/flat run.
Once that shared construction has been completed,
it gives, on this very frame,
\[
 \mathrm{FinPt}(\Phi,3),\qquad
 \mathrm{AutPt}(\Phi,3,\epsilon_\ell^{-1}),\qquad
 \mathrm{AutPt}(\Phi,2,k_2).
 \tag{G10}
\]
The two actual distinguished auxiliary places remove
each other's omitted-place qualification.
This is the precise use of the common theorem;
(G10) is not inferred merely from an Eisenstein
congruence or a regular algebraic realization of the seed.

\subsubsection{From the shared output to the Gaussian host endpoint}
We record the final implication so that, after (G10),
no further Gaussian automorphy premise remains.
The quotient map in \(\mathrm{GalLink}\) is by the
closed normal closure of the inertia elements on which
\(\eta\) and \(\bar\Psi\) are both trivial.
An actual marked Gaussian point therefore restricts
to the exact ordinary rank-three representation
of \(\Phi\), with inertial tuple
\((\Psi,1,\Psi^{-1})\).
The source's \(\mathrm{exists\_nrep3}\) is this
restriction and ordinary flag, with its stable lattice;
it does not posit a new representation.

First use \(\mathrm{FinPt}(\Phi,3)\) in (G10) to obtain
\(\mathrm{FinEnd}\) on the open group \(G_{L_r}\).
Characteristic-\(\ell\) fixed-inertia points then have
the residual characteristic polynomials there.
Lemma~\ref{lem:sr-constant-character-triangular}
gives upper or lower triangularity in the actual marked
basis and fixes each diagonal to its constant character.
The balanced-monomial argument in
Lemma~\ref{lem:sr-graded-finite-part-criterion},
applied exactly in
Proposition~\ref{prop:sr-host-finite-part-from-finend},
therefore makes the weight-zero universal part finite
over the coefficient ring.
These proofs reconstruct the needed
\srcid{ConstOfChar}--\srcid{UniversalDiagonalOfTriangularPoints}--%
\srcid{FinCrit} chain and are valid in characteristic three.
It is exactly the required specialization of the
finiteness clause of \(\mathrm{ThmS6Core\_s}\).
Every resulting characteristic-zero universal point
therefore admits the finite integral DVR model and
normalization used by the ordinary-point argument.
This finiteness is obtained before applying the
strong-open assertion to such a point.

For a type-I characteristic-zero host point,
Section~\ref{sec:spectrum-open-irreducibility} proves
irreducibility on every open subgroup of \(G_K\).
Thus its restriction to \(G_{L_r}\) is irreducible.
The rank-three assertion in (G10) gives actual
regular algebraic cuspidal automorphy of
\[
                             \widehat\rho|_{G_{L_r}},
                         \qquad \widehat\rho=\rho\epsilon_\ell^{-1}.
\]
Use exact cyclic descent
(Lemma~\ref{lem:sr-cyclic-exact-descent}) along
the abelian CM extension \(L_r/K\).
The finite character correction in that lemma is
included, so the descended representation is attached
exactly to \(\widehat\rho\).

The semidirect unitary law and the normalized conjugation
matrix give
\[
              \rho(g)^{\mathsf t}J_3\rho(cgc^{-1})=J_3
                           \quad(g\in G_K).
\]
Thus \(\widehat\rho\) has the exact conjugate pairing
with multiplier \(\epsilon_\ell^{-2}\).
Chebotarev and strong multiplicity one give genuine
conjugate self-duality of its rank-three automorphic
realization.
Full local--global compatibility therefore applies.
Its ordinary, finite-inertia, and full local group
conditions give the actual local window predicate,
as in Proposition~\ref{prop:sr-actual-wit};
the scalar normalization cancels from every ratio.
At \(\ell=3,7\) one may also use directly
Lemma~\ref{lem:sr-gaussian-local-ratios} and the
full comparison: there is no cyclotomic edge,
monodromy is zero, and the fixed inertial type is
unramified in Weil--Deligne normalization.
This yields the asserted crystalline condition.
No comparison of semisimplified Weil characters alone
is used to identify monodromy.

The type-I points have the automorphy just descended,
and their strong-open irreducibility was proved
independently in the cited ordinary-point argument.
Proposition~\ref{prop:sr-cyclotomic-restriction} supplies
the infinite-cyclotomic-kernel condition.
Together these give all three clauses of
\(\mathrm{BPrime\_s}\) at the actual Gaussian host,
once the explicitly stated common support input (G10)
is supplied.

This locates the remaining common dependency exactly.
The Gaussian prime-\(3\) seed, its character normalization,
its actual module insertion, the finite-unramified
real quadratic base, local frame witnesses and the
descent/local-comparison implication have been proved
here. Completion of the shared sharp/flat propagation
construction supplies (G10); none of the predicates
\(\mathrm{BPrime\_s.aut}\) or
\(\mathrm{ThmS6Core\_s}\) is used as a substitute for it.

% Frozen source:
% DnEndsTop.ends0_alt -> DnEnds.ends_u_alt -> SymTop.linkedEnd_alt
% -> LinkedEndE.linked_endomorphism_dense -> TwoRuns.twoRuns_alt.
% Seed2Alt: SdStmt, BirthTop, BirthEG, IsSeed.
% Host interface: DnHost, AutEnd3Alt, SymmetricCubeTheoremOfEndomorphisms.
% Actual local conditions: FrobeniusSmoothVectors/
% FrobeniusSplitPrimeDatum.Gen.g4,g5 and UFacts.mu,f3.


\subsection{The actual generic rank-three endpoint after propagation}
\label{subsec:sr-generic-three-endpoint}
We spell out the endpoint over the fixed field
\(K=\mathbb Q(\sqrt{-3})\).
Let \(a\ge1\), let \(\ell\) and \(\psi\) satisfy the
good-data hypotheses of Section~\ref{sec:single-character},
and put \(\Psi=\epsilon_\ell^a\psi\).
Let \(R\) be the actual marked rank-three point, let
\(\rho=R|_{G_K}\), and put
\(\widehat\rho=\rho\epsilon_\ell^{-1}\).
Assume the common sharp/flat construction has supplied,
on the actual frame and actual ordinary function modules,
\[
 \mathrm{FinPt}(\Phi,3),\qquad
 \mathrm{AutPt}(\Phi,3,\epsilon_\ell^{-1}).
 \tag{GE1}
\]
These are conclusions of the propagation construction,
rather than hypotheses about an unrelated automorphic
representation.

\begin{srproposition}\label{prop:sr-generic-three-exact-endpoint}
Under (GE1), the weight-zero part of the marked host
algebra is finite over the coefficient DVR.
Every type-I characteristic-zero point gives an actual
regular algebraic conjugate self-dual cuspidal representation
over \(K\) attached exactly to \(\widehat\rho\).
It satisfies the full local predicate \(\mathrm{WitA}\).
Moreover \(\widehat\rho\) is irreducible on
\(G_{K(\mu_{\ell^\infty})}\).
\end{srproposition}
\begin{proof}
The map from \(G_{\mathbb Q}\) to the frame group is the
actual Galois link.
Its restriction carries the marked point, its ordinary
flag and its finite inertia characters to the normalized
rank-three frame point.
Thus the first assertion of (GE1) fixes the characteristic
polynomials on the open subgroup \(G_L\subset G_K\).
For characteristic-\(\ell\) points,
Proposition~\ref{prop:sr-host-finite-part-from-finend}
proves the required finite weight-zero algebra.
Its proof uses the actual marked generators, so applies
to the generic host as well as the Gaussian host.
The integral normalization of an ordinary characteristic-zero
point can now be made as in
Section~\ref{sec:spectrum-deformation}.

For a type-I point the residual inertia separator is
\[
 \exists u\in I_\ell:\quad \overline\Psi(u)^2\ne1.
 \tag{GE2}
\]
Indeed the good-data condition \(\psi_\ell^2=1\)
gives
\[
 \overline\Psi|_{I_\ell}
      =\overline\epsilon_\ell^{\,a+b},
       \qquad b\in\{0,(\ell-1)/2\},
 \quad
 \overline\Psi^{\,2}|_{I_\ell}
      =\overline\epsilon_\ell^{\,2a}.
\]
This is the local quadratic character calculation in the
actual principal packet.
The good-data condition \(\ell-1\nmid2a\) therefore
gives (GE2), including when \(\ell\mid N\).
Also \(I_\ell\subset G_K\), since \(\ell\ne3\).
Section~\ref{sec:spectrum-open-irreducibility},
with exactly this separator and the ordinary flag,
proves irreducibility on every open subgroup of \(G_K\).
In particular \(\widehat\rho|_{G_L}\) is irreducible.
The second assertion of (GE1) therefore gives actual
regular algebraic cuspidal automorphy of
\(\widehat\rho|_{G_L}\).
The twist in (GE1) is already \(\epsilon_\ell^{-1}\):
the geometric-Frobenius convention gives
\[
 r_{\ell,\iota}(\operatorname{Ad}\pi_f)
       =\operatorname{ad}^0(r_{\ell,\iota}(\pi_f))
                                      \epsilon_\ell^{-1},
\]
which is the normalization of the rank-three seed.

The extension \(L/K\) is finite abelian and CM.
Apply Lemma~\ref{lem:sr-cyclic-exact-descent}
successively along its cyclic composition series.
The invariant representation at a descent step is the
restriction of the original \(\widehat\rho\), and its
restriction to \(G_L\) remains irreducible.
Consequently the lemma includes its finite-character
correction and gives a cusp representation attached
exactly to \(\widehat\rho\), with no undetermined
finite twist left at the endpoint.

The normalized semidirect-product law gives
\[
 \rho(cgc^{-1})=J_3\rho(g)^{-\mathsf t}J_3^{-1}.
\]
Since \(\epsilon_\ell(cgc^{-1})=\epsilon_\ell(g)\),
\[
 \widehat\rho(cgc^{-1})
       =\epsilon_\ell(g)^{-2}
                   J_3\widehat\rho(g)^{-\mathsf t}J_3^{-1}.
 \tag{GE3}
\]
The exponent is \(1-3=-2\), exactly the prescribed
rank-three Galois normalization.
The pairing is symmetric and its real multiplier is
\(\epsilon_\ell^{-2}\).
Chebotarev and strong multiplicity one therefore make
the descended regular algebraic cusp representation
genuinely conjugate self-dual.
Full local--global compatibility applies to that
actual polarized representation, including monodromy
at places above \(\ell\).

Here are the local clauses; they use no comparison of
semisimplified Weil characters to identify monodromy.
Away from \(\ell\), the actual ramification condition
kills inertia intersected with the kernel of the fixed
residual torus tuple and the quadratic sign.
This kernel contains the global commutator subgroup.
The resulting inertia image is finite abelian and its
isotypic subspaces are decomposition-stable.
Thus its Frobenius-semisimplified Weil representation
is a sum of characters and has \(N=0\).
At a place above \(\ell\), the ordinary flag has parallel
cyclotomic exponents
\[
                         (a-1,-1,-a-1)
 \tag{GE4}
\]
after the scalar twist.
The de Rham assertion comes from the actual automorphic
realization.
If the flag satisfies
\(\kappa_i\ne\kappa_j\epsilon_\ell\) for all \(i,j\),
Lemma~\ref{lem:sr-parallel-wd-characters} implies \(N=0\).

A single real norm twist makes the global cusp representation
unitary. At every place with this rank-one parameter
and \(N=0\), its generic unitary local representation
has all its inducing characters in the strict window
\[
                  q_v^{\,t-1/2}<|\xi(\varpi_v)|
                                         <q_v^{\,t+1/2}
 \tag{GE5}
\]
for the same real \(t\).
This is the character-parameter instance of the
unitary classification used in
Proposition~\ref{prop:sr-actual-wit}.
An equality
\((\kappa_i/\kappa_j)^m=\epsilon_\ell^m\)
would give a Frobenius absolute-value ratio \(q_v\)
in one order, contrary to (GE5).
Thus \(\mathrm{Unl1}\) implies \(\mathrm{PotUnl}\),
and the same \(t\) gives the window at every place
away from \(\ell\), and at every \(\ell\)-adic place
satisfying \(\mathrm{PotUnl}\).
Any other attached cusp representation has the same
unramified Hecke factors and equals this one by strong
multiplicity one.
This proves every quantifier in \(\mathrm{WitA}\),
without assuming \(\mathrm{Unl1}\) at a place where
it has not been established.

Finally the open-subgroup irreducibility already proved,
and Proposition~\ref{prop:sr-cyclotomic-restriction},
give irreducibility on \(G_{K(\mu_{\ell^\infty})}\).
The scalar inverse-cyclotomic twist changes no invariant
subspace.
\end{proof}

\paragraph{The separate original-level clause.}
The preceding proposition gives the actual
\(\mathrm{AutEnd3Alt}\) and \(\mathrm{IsAutCsdA}\)
endpoint, including its complete local window.
For \(a\ge2\), the differences of the three exponents
in (GE4) have absolute values \(a,a,2a\), none equal
to one. Thus \(\mathrm{Unl1}\) is automatic, and
(GE5) applies at every \(\ell\)-adic place.
For \(a=1\), this Hodge-number argument cannot be used:
nearest differences are one.

The source's stronger counting class
\(\mathrm{Adm3U_s}\) also requires
\[
 \mathrm{Th.Unl}\bigl(\ell,N,(a,0,-a),\rho\bigr).
 \tag{GE6}
\]
In the archive this is supplied by
\(\mathrm{UnlEndLogNorm}\).
An alternative sufficient producer in the present
argument is the actual local component assertion
\[
 \begin{gathered}
 I_Y\subset I_{\mathrm{cr}}\subset I_t,\qquad
 I_t\subset\ker x
 \quad\Longrightarrow\\[-2mm]
 \rho_x|_{\ker\eta\cap\ker\tau}
                         \text{ is crystalline}.
 \end{gathered}
 \tag{GE7}
\]
Here \(I_Y\) is the ordinary ideal saturated by
\(\ell\,dd_0dd_1\), \(I_{\mathrm{cr}}=\sqrt{I_Y}\),
and \(I_t\) is the selected component through the torus
point, all in the actual local universal ring.
Once (GE7), including the factorization of every
selected global point through \(I_t\), is proved,
its crystalline monodromy is zero.
Full local--global compatibility then yields (GE5)
without assuming \(\mathrm{Unl1}\).
That strict window prohibits
\(\kappa_i=\kappa_j\epsilon_\ell\) itself, proving
(GE6) also when \(a=1\).
The specified finite inertial type gives the original
local level.
When \(\ell\nmid N\), that type is trivial,
so the attached representation is crystalline and its
local automorphic component is unramified.
When \(\ell\mid N\), one retains the prescribed
finite type \((\psi_\ell,1,\psi_\ell^{-1})\):
the representation is crystalline after restriction to
\(\ker\eta\cap\ker\tau\), hence potentially crystalline
over the original local field, and its finite conductor
is accounted for by the existing \(\ell\)-part of \(N\).
No unramifiedness assertion is made in that case.
This is precisely what prevents an unwanted factor
of the varying coefficient prime from entering the
common integer's fixed spectrum.

Neither ordinarity alone nor the name
\(I_{\mathrm{cr}}\) proves (GE7).
The theorem
\srcid{CrystallineOrdinaryIdeals/RtPoints.lean}
uses the separate normal-form result
\srcid{ideal\_cuts\_crystalline\_locus};
its proof and the actual selected-chart factorization
are the remaining local producers for this alternative.
They are kept distinct here from the support output
(GE1).
